% ==================/// PHASES ///==================
\section{Phases }
% ==================/// PHASE 1 ///==================
\begin{frame}{Phase 1: TF setup and perception}
    \begin{columns}[T]

        % LEFT: one step visible at a time
        \begin{column}{0.6\textwidth}
            {\small

            \only<1>{
                \textbf{Step 1 — Put tags on the map.}\\[0.2cm]
                We save where each AprilTag sits:
                \[
                    T_{\text{world}}^{\text{tag}}
                \]
                This stays fixed for every next step.
            }

            \only<2>{
                \textbf{Step 2 — Drive the robot once.}\\[0.2cm]
                The simulator gives the robot pose while it follows a smooth curve:
                \[
                    T_{\text{world}}^{\text{camera\_link}}(t)
                \]
            }

            \only<3>{
                \textbf{Step 3 — Use the real camera frame.}\\[0.2cm]
                We switch to the optical frame (how the camera really looks):
                \[
                    T_{\text{world}}^{\text{camera\_optical}}(t)
                    =
                    T_{\text{world}}^{\text{camera\_link}}(t)
                    \cdot
                    T_{\text{camera\_link}}^{\text{camera\_optical}}
                \]
            }

            \only<4>{
                \textbf{Step 4 — Compute clean tag detections.}\\[0.2cm]
                From the map and pose we get, for each tag:
                \[
                    T_{\text{camera\_optical}}^{\text{tag}}(t)
                \]
                This is the clean data that Phase 2 will dirty with noise.
            }

            }
        \end{column}

        % RIGHT: image appears from step 2 onwards and stays
        \begin{column}{0.35\textwidth}
            \centering
            \only<2->{
                \includegraphics[width=\textwidth]{assets/TF.png}
                \vspace{0.15cm}

                {\tiny\textit{Fig. 4 — World, robot and camera frames}}
            }
        \end{column}

    \end{columns}
\end{frame}



% ==================/// PHASE 2 ///==================
\begin{frame}{Phase 2: Simulation and dataset generation}
	\begin{columns}[T]
		
		% LEFT COLUMN
		\begin{column}{0.58\textwidth}
			\small
			
			\only<1>{
				\textbf{Step 1 — Simulate a clean path.}\\[0.2cm]
				We make a smooth curve and log the true pose:\\[0.1cm]
				\[
				T_{\text{world}}^{\text{camera\_optical}}(t)
				\]
			}
			
			\only<2>{
				\textbf{Step 2 — Compute perfect tag reads.}\\[0.2cm]
				For every tag and time we compute the clean detection:\\[0.1cm]
				\[
				T_{\text{camera\_optical}}^{\text{tag}}(t)
				\]
			}
			
			\only<3>{
				\textbf{Step 3 — Add noise like real life.}\\[0.2cm]
				We add Gaussian noise so readings look imperfect, like sensors do.
			}
			
			\only<4->{
				\textbf{Step 4 — Save the dataset.}\\[0.2cm]
				For each time we save:\\[0.1cm]
				- the true pose $T_{\text{world}}^{\text{camera\_optical}}(t)$\\
				- the noisy tag detections $T_{\text{camera\_optical}}^{\text{tag}}(t)$\\[0.2cm]
				Phase 3 will test its optimizer on this file.
			}
			
		\end{column}
		
		% RIGHT COLUMN — KEEP BIG IMAGES
		\begin{column}{0.40\textwidth}
			\centering
			
			\begin{overlayarea}{\textwidth}{0.60\textheight}
				
				% CLEAN IMAGE (large)
				\only<1>{
					\includegraphics[width=\textwidth]{assets/s_curve_dataset_large_plot_big.png}\\
					{\tiny\textit{Fig. 5 — Ground truth sinusoid trajectory}}
				}
				
				% NOISY IMAGE (large)
				\only<2>{
					\includegraphics[width=\textwidth]{assets/s_curve_dataset_large_noisy_plot_big.png}\\
					{\tiny\textit{Fig. 6 — Noisy measurements}}
				}
				
				% TWO BIG IMAGES STACKED (one above the other)
				\only<3->{
					\includegraphics[width=0.6\textwidth]{assets/s_curve_dataset_large_plot_big.png}\\[0.2cm]
					\includegraphics[width=0.6\textwidth]{assets/s_curve_dataset_large_noisy_plot_big.png}
					
					{\tiny\textit{Clean (top) and noisy (bottom) simulated data}}
				}
				
			\end{overlayarea}
		\end{column}
		
	\end{columns}
\end{frame}




% ==================/// PHASE 3 ///==================
\begin{frame}{Phase 3: Trajectory Optimization \& Recovery}
    \begin{columns}
        \begin{column}{0.6\textwidth}
            \textbf{Goal: Fix drift with a simple optimizer}
            \vspace{0.5cm}
            \begin{itemize}
                \item<1-> \textbf{Path recovery:}  
                Align the noisy run with the taught path.

                \item<2-> \textbf{Pose graph:}  
                Reduce error between waypoints and simple constraints.

                \item<3-> \textbf{Drift fix:}  
                Auto-correct rotation and position drift.
            \end{itemize}
        \end{column}
        
        \begin{column}{0.45\textwidth}
            \centering
            % A imagem aqui deve ser o gráfico de otimização gerado acima
            \only<2->{
                \includegraphics[trim={5cm 2cm 5cm 2cm}, clip, width=\textwidth]{assets/revovery_path.png}
                \vspace{0.2cm}
                {\tiny\textit{Fig. 7 — Recovery of the optimal path via algorithm}}
            }
        \end{column}
    \end{columns}
\end{frame}


% ==================/// PHASE 4 ///==================
\begin{frame}{Phase 4: Robot implementation}
    \begin{columns}[T] % [T] Alinha o conteúdo pelo topo
        
        % --- Coluna da Esquerda (Texto) ---
        \begin{column}{0.6\textwidth}
            \textbf{Goal: Move from PC tests to the LIMOS robot}
            \vspace{0.5cm}
            \begin{itemize}
                \setlength\itemsep{1em} % Espaçamento entre itens
                
                \item<1-> \textbf{Onboard integration:}  
                Put the perception code on the robot (ROS~2).

                \item<3-> \textbf{Camera calibration:}  
                Calibrate the camera so distances look right.

                \item<5-> \textbf{First trajectory recording:}  
                Record one path using wheels plus AprilTag pose.
            \end{itemize}
        \end{column}
        
        % --- Coluna da Direita (Imagens) ---
        \begin{column}{0.38\textwidth}
            \centering
            % overlayarea reserva o espaço para evitar que o slide pule na troca
            \begin{overlayarea}{\textwidth}{0.6\textheight}
                \centering
                
                % IMAGEM 1: Aparece no slide 2 e SOME no final do slide 3
                \only<2>{
                    \includegraphics[width=\textwidth, keepaspectratio]{assets/limo.jpeg}
                    \vspace{0cm}
                    \\{\tiny\textit{Fig. 8 — LIMOS platform}}
                }
                
                % IMAGEM 2: Aparece no slide 4 e fica até o fim
                % (Coloquei o mesmo nome de arquivo, mas idealmente seria outra foto)
                \only<4->{
                    \includegraphics[trim={5cm 3cm 5cm 2cm}, clip, width=\textwidth]{assets/camera_calibration.jpg}
                    \vspace{0cm}
                    \\{\tiny\textit{Fig. 9 — Calibration}}
                }
            \end{overlayarea}
        \end{column}
        
    \end{columns}
\end{frame}



% ==================/// PHASE 5 ///==================
\begin{frame}{Phase 5: Path recovery algorithm}
    \begin{columns}
        \begin{column}{0.6\textwidth}
            \textbf{Goal: Recover when the robot leaves the taught path}
            \vspace{0.5cm}
            \begin{itemize}
                \item<1-> \textbf{Recovery logic:}  
                Simple rules to steer back to the recorded path.

                \item<2-> \textbf{Closest waypoint search:}  
                Use a quick KD-Tree to find the nearest waypoint.

                \item<4-> \textbf{Safety validation:}  
                Test recovery moves in sim to avoid bad turns.
            \end{itemize}
        \end{column}
        
        \begin{column}{0.50\textwidth}
            \centering
            \only<3->{
                \includegraphics[trim={5cm 4cm 5cm 3cm}, clip, width=\textwidth]{assets/return_path.png}

                \vspace{0.2cm}
                {\tiny\textit{Fig. 10 — Returning to the original path after deviation}}
            }
        \end{column}
    \end{columns}
\end{frame}



% ==================/// PHASE 6 ///==================
\begin{frame}{Phase 6: System integration}
    \begin{columns}
        \begin{column}{0.6\textwidth}
            \textbf{Goal: Demonstrate full teach and repeat autonomy}
            \vspace{0.5cm}
            \begin{itemize}
                \item<1-> \textbf{Controller integration:}  
                Mix Pure Pursuit with the perception and pose stack.

                \item<2-> \textbf{Autonomous replay:}  
                Let the robot drive the recorded path on its own.

                \item<4-> \textbf{Final system:}  
                A teach \& repeat pipeline anyone can run without coding.
            \end{itemize}
        \end{column}
        
        \begin{column}{0.45\textwidth}
            \centering
            \only<3->{
                \includegraphics[trim={1cm 5cm 1cm 1cm}, clip, width=\textwidth]{assets/robot_transit.png}

                \vspace{0.2cm}
                {\tiny\textit{Fig. 11 — Fully integrated teach and repeat behavior}}
            }
        \end{column}
    \end{columns}
\end{frame}
